\chapter*{Tóm tắt}
\label{summary}

Nhiều sinh viên của trường Đại học Khoa học Tự nhiên – ĐHQG TPHCM đang học tập và làm việc trong tình trạng không có định hướng cho bản thân. Hệ thống hỗ trợ cung cấp thông tin học tập ra đời sẽ giúp sinh viên có góc nhìn tổng quan hơn trong con đường học vấn. Hệ thống sẽ đưa ra được các thông tin học tập theo chuyên ngành hiện tại, mong muốn sau này của sinh viên, gợi ý lộ trình học tập rõ ràng, chi tiết để đạt được hiệu quả học tập tốt nhất trong quá trình học. Nghiên cứu này áp dụng Knowledge Graph (KG) để lưu trữ thông tin học vụ, kết hợp AI Agents để chuyển hóa tri thức thành đồ thị và hỗ trợ tìm kiếm, phân tích lộ trình học tập tối ưu cho người dùng. Về mặt phương pháp, nghiên cứu cải tiến quy trình trích xuất đồ thị tri thức bằng mô hình biểu diễn bộ tứ (Quadruple) kết hợp với luồng tối ưu hóa prompt (Prompt Flow) dành riêng cho các tài liệu chuyên biệt. Hệ thống áp dụng kỹ thuật phân mảnh văn bản dựa trên vector ngữ nghĩa (Chunk Vector Section) để duy trì tính toàn vẹn của ngữ cảnh, đồng thời tích hợp cơ chế tự động giải quyết và đồng nhất các từ viết tắt liên quan. Để đánh giá toàn diện hệ thống, nghiên cứu tiến hành đánh giá trên hai phương diện: chất lượng Đồ thị tri thức (KG) và hiệu năng của hệ thống RAG. Quá trình này sử dụng phương pháp giám khảo tự động (LLM-as-a-Judge) kết hợp với khung đánh giá chuẩn RAGAS. Kết quả thực nghiệm trên KG cho thấy kiến trúc đề xuất đã cải thiện đáng kể chất lượng liên kết, khắc phục thành công và làm giảm mạnh tỷ lệ nút mồ côi (dangling nodes). Bên cạnh đó, các tiêu chí đánh giá RAG về độ liên quan và độ phủ ngữ cảnh cũng đạt mức cao, chứng minh khả năng cung cấp thông tin tư vấn học vụ chính xác và cá nhân hóa trên nền tảng ứng dụng.
